\chapter{Adição: primeiros passos}\label{ch:g1:addition}

Três bolinhas vermelhas e duas bolinhas azuis: juntas, quantas são?
Juntar e contar o total é \emph{somar} --- a primeira operação, com
os seus dois sinais famosos $+$ e $=$.

\section{Juntar}

\begin{definition}[Adição]\label{def:g1:addition:def}
\emph{Somar} é juntar e contar tudo. Escrevemos
\[
3 + 2 = 5
\]
e lemos: ``três mais dois é igual a cinco''. O resultado de uma
adição chama-se \emph{soma}\index{soma}.
\end{definition}

\begin{omfigure}
\begin{tikzpicture}[scale=0.6]
  \foreach \i in {0,1,2} \fill[omThm] (\i*1.0,0) circle (0.3);
  \node[font=\large] at (3.0,0) {$+$};
  \foreach \i in {0,1} \fill[omDef] (3.8+\i*1.0,0) circle (0.3);
  \node[font=\large] at (5.8,0) {$=$};
  \foreach \i in {0,1,2} \fill[omThm] (6.6+\i*1.0,0) circle (0.3);
  \foreach \i in {0,1} \fill[omDef] (9.6+\i*1.0,0) circle (0.3);
  \node[below, font=\small] at (5.6,-0.7)
    {$3 + 2 = 5$: três e dois fazem cinco};
\end{tikzpicture}

{\small Uma adição em desenho: os dois grupos juntos formam um único
grupo de cinco.}
\end{omfigure}

\begin{example}[A ordem não importa]\label{ex:g1:addition:order}
$2 + 3$ e $3 + 2$ juntam as mesmas bolinhas: as duas contas dão $5$.
Para calcular $2 + 9$, é mais rápido começar pelo número maior: diga
``nove'' e conte mais dois --- ``dez, onze''. Logo $2 + 9 = 11$.
\end{example}

\section{Os amigos do dez}

\begin{example}[Complementos para 10]\label{ex:g1:addition:friends}
Dois números que juntos fazem $10$ são ``amigos do dez'':
\[
1 + 9, \quad 2 + 8, \quad 3 + 7, \quad 4 + 6, \quad 5 + 5 .
\]
Os dez dedos mostram todos eles: levante $3$ dedos e os $7$ dobrados
são o amigo. Saber os amigos do dez de cor deixa muitas adições
rápidas.
\end{example}

\begin{omfigure}
\begin{tikzpicture}[scale=0.72]
  % quadro de dez: 2 fileiras de 5 casas
  \draw[thick] (0,0) grid[step=1] (5,2);
  \draw[very thick] (0,0) rectangle (5,2);
  \foreach \i in {0,1,2} \fill[omDef] (\i+0.5,1.5) circle (0.3);
  \foreach \i in {3,4} \draw[omThm, thick] (\i+0.5,1.5) circle (0.3);
  \foreach \i in {0,...,4} \draw[omThm, thick] (\i+0.5,0.5) circle (0.3);
  \node[right, font=\small, align=left] at (5.6,1.0)
    {um \emph{quadro do dez}: $10$ casas.\\
     $3$ cheias, $7$ vazias: $3 + 7 = 10$};
\end{tikzpicture}

{\small O quadro do dez mostra num relance cada par de amigos do dez:
o que está cheio e o que falta sempre fazem $10$ juntos.}
\end{omfigure}

\begin{method}[Passar pelo dez]\label{met:g1:addition:crossten}
Para calcular $8 + 5$:
\begin{enumerate}
  \item tire do $5$ o que falta ao $8$ para chegar a $10$: são $2$
        (o amigo do dez do $8$);
  \item $8 + 2 = 10$, e ainda sobram $3$ do $5$;
  \item $10 + 3 = 13$. Logo $8 + 5 = 13$.
\end{enumerate}
Passar pelo $10$ transforma uma \omterm{def:g1:addition:def}{soma} difícil em duas \omterm{def:g1:addition:def}{somas} fáceis.
\end{method}

\begin{omfigure}
\begin{tikzpicture}[scale=0.62]
  % quadro de dez com 8 azuis e 2 vermelhas (emprestadas do 5)
  \draw[thick] (0,0) grid[step=1] (5,2);
  \draw[very thick] (0,0) rectangle (5,2);
  \foreach \i in {0,...,4} \fill[omDef] (\i+0.5,1.5) circle (0.3);
  \foreach \i in {0,1,2} \fill[omDef] (\i+0.5,0.5) circle (0.3);
  \foreach \i in {3,4} \fill[omThm] (\i+0.5,0.5) circle (0.3);
  % 3 vermelhas ficam fora do quadro
  \foreach \i in {0,1,2} \fill[omThm] (6.2+\i,1.0) circle (0.3);
  \node[below, font=\small] at (2.5,-0.25) {o quadro está cheio: $10$};
  \node[below, font=\small] at (7.2,0.55) {sobram $3$};
\end{tikzpicture}

{\small $8 + 5$ no quadro do dez: as $8$ fichas azuis tomam $2$ das
$5$ vermelhas para encher o quadro --- uma dezena cheia e mais $3$,
ou seja $13$.}
\end{omfigure}

\begin{omfigure}
\begin{tikzpicture}[scale=0.6]
  \draw[-Stealth] (-0.4,0) -- (15.4,0);
  \foreach \i in {0,...,15} \draw (\i,-0.1) -- (\i,0.1);
  \foreach \i in {0,5,8,10,13,15}
    \node[below, font=\footnotesize] at (\i,-0.12) {\i};
  \draw[-Stealth, omDef, thick] (8,0.3) .. controls (9,0.95) ..
    (10,0.3);
  \node[omDef, font=\small] at (9,1.05) {$+2$};
  \draw[-Stealth, omThm, thick] (10,0.3) .. controls (11.5,1.15) ..
    (13,0.3);
  \node[omThm, font=\small] at (11.5,1.2) {$+3$};
\end{tikzpicture}

{\small $8 + 5$ na reta numérica: pule primeiro até o $10$, depois o
resto.}
\end{omfigure}

\begin{example}[Dobros]\label{ex:g1:addition:doubles}
Vale a pena saber os dobros de cor:
\[
1+1=2, \quad 2+2=4, \quad 3+3=6, \quad 4+4=8, \quad 5+5=10,
\]
\[
6+6=12, \quad 7+7=14, \quad 8+8=16, \quad 9+9=18, \quad 10+10=20 .
\]
Os quase-dobros saem de graça: $6 + 7$ é um a mais que $6 + 6$, ou
seja $13$.
\end{example}

\section{Probleminhas}

\begin{example}\label{ex:g1:addition:problem}
``Lia tem $7$ adesivos. Ganha mais $6$. Com quantos fica?'' Ganhar
mais é somar: $7 + 6 = 13$ (passando pelo dez: $7 + 3 = 10$, depois
$+3$). Lia fica com $13$ adesivos. Termine sempre com uma frase de
resposta!
\end{example}

\section{Exercícios}

\begin{exercise}[$\star$]\label{exo:g1:addition:1}
Desenhe as bolinhas e calcule: $4 + 3$; \quad $5 + 2$; \quad
$6 + 4$.
\end{exercise}

\begin{exercise}[$\star$]\label{exo:g1:addition:2}
Calcule começando pelo número maior: $3 + 8$; \quad
$2 + 12$; \quad $4 + 9$.
\end{exercise}

\begin{exercise}[$\star$]\label{exo:g1:addition:3}
Dê o amigo do dez de: $7$; \ $4$; \ $9$; \ $5$; \ $1$.
\end{exercise}

\begin{exercise}[$\star$]\label{exo:g1:addition:4}
Calcule passando pelo dez (\cref{met:g1:addition:crossten}):
$9 + 4$; \quad $8 + 7$; \quad $6 + 5$.
\end{exercise}

\begin{exercise}[$\star$]\label{exo:g1:addition:5}
Recite os dobros e depois use-os para calcular: $7 + 8$; \quad
$5 + 6$; \quad $9 + 8$.
\end{exercise}

\begin{exercise}[$\star$]\label{exo:g1:addition:6}
Complete os buracos:
\[
6 + \;?\; = 10, \qquad
\;?\; + 5 = 12, \qquad
10 + \;?\; = 17 .
\]
\end{exercise}

\begin{exercise}[$\star$]\label{exo:g1:addition:7}
``No ônibus há $8$ crianças; $5$ adultos sobem. Quantas pessoas há
no ônibus?'' Escreva a adição e a frase de resposta.
\end{exercise}

\begin{exercise}[$\star$]\label{exo:g1:addition:8}
Tomás tem $6$ carrinhos vermelhos e $7$ azuis. Mia tem $10$
carrinhos. Quem tem mais carrinhos?
\end{exercise}

\begin{exercise}[$\star\star$]\label{exo:g1:addition:9}
Três dados mostram $4$, $6$ e $5$. Qual é o total? (Some primeiro
dois dados --- escolha os dois que são amigos do dez!)
\end{exercise}
